\chapter{Pascalian Roots of Unity}

\vspace{-\baselineskip}
\lettrine{M}{arjorie} Worpitsky’s investigations into Pascalian roots of unity had uncovered surprising connections to recursive structures and polynomial symmetries. 

\begin{figure}[H]
\includegraphics[width=1.0\linewidth]{Illustrations/vign-samLau2.png}
\caption{Sam Lau with et al number one}
\end{figure}

But when Sam Lau et al introduced a complex Fibonacci matrix into their discussions, Marjorie saw a new dimension unfolding—one that tied the twisting symmetries of complex numbers to the structural elegance of her Pascalian framework. 


\lettrine{S}am presented to his other femme et al the following matrix:
\[
A = 
\begin{pmatrix}
1 & i & 0 \\
i & 1 & i \\
0 & i & 1
\end{pmatrix}.
\]

It was this matrix that \(\mathbb{C}\)omplexified the Fibonacci structure, broadening potential symmetries with this golden recursion. He remarked that it was often the case in maths that in broadening one's viewpoint through some functor operation initial clarity is lost but some greater insight is usually gleaned later when the functor is reversed. 

\begin{wrapfigure}{r}{0.7\textwidth}
    \centering
    \includegraphics[width=\linewidth]{Illustrations/vign-sam.png}
    \caption{Sam Lau et al number two}
    \label{fig:Femme}
\end{wrapfigure}
 
 Sam was in communication with Abu Rashid who had also analyzed its properties some years ago. Together with their group of postgrads they began uncovering connections to the works of Demi and Marj. 
 
 Fifty years hence Abu hoped to be referred to kindly-in circles that might have meant something to them - as the <<AbuLau>> group. Sam merely hoped that they might be looked upon with some fondness, as a product of their time.
 
 Abu was particularly keen to keep the communication open and receptive with Marj especially when she began to blossom in the twilight hours which was happily when he was at the peak of his abilities. He calculated - it was probably his best calculation of the month - that she may well have keeping some cards close to her chest and he was hoping his insights and charm may persuade her to reveal them.  

\paragraph{Hermitian Symmetry:} The matrix \(A\) is Hermitian, satisfying:
\[
A^\dagger = A,
\]
where \(A^\dagger\) is the conjugate transpose of \(A\). This property, common in quantum mechanics, suggested that \(A\) might correspond to observable symmetries in a complex geometric or algebraic space.

\paragraph{Determinant:} The determinant of \(A\) measures its scaling effect in complex space. Calculating it:
\[
\det(A) = \begin{vmatrix}
1 & i & 0 \\
i & 1 & i \\
0 & i & 1
\end{vmatrix}.
\]
Expanding along the first row:
\[
\det(A) = 1 \cdot \begin{vmatrix}
1 & i \\
i & 1
\end{vmatrix}
- i \cdot \begin{vmatrix}
i & i \\
0 & 1
\end{vmatrix}
+ 0 \cdot \begin{vmatrix}
i & 1 \\
0 & i
\end{vmatrix}.
\]
Evaluating the minors:
\[
\begin{vmatrix}
1 & i \\
i & 1
\end{vmatrix} = 1 \cdot 1 - i \cdot i = 1 - (-1) = 2,
\]
\[
\begin{vmatrix}
i & i \\
0 & 1
\end{vmatrix} = i \cdot 1 - i \cdot 0 = i.
\]
Substituting back:
\[
\det(A) = 1 \cdot 2 - i \cdot i = 2 - (-1) = 3.
\]
The determinant, a real number, hinted at a balanced structure despite the matrix’s complex entries.

\subsection*{Generating a Complex Fibonacci Sequence}

The matrix \(A\) can produce a Fibonacci-like sequence by repeated application to an initial vector, \(v_0 = \begin{pmatrix} 1 \\ 0 \\ 0 \end{pmatrix}\):
\[
v_{n+1} = A v_n.
\]
Computing the first few terms:

\paragraph{First Iteration:}
\[
v_1 = A v_0 = 
\begin{pmatrix}
1 & i & 0 \\
i & 1 & i \\
0 & i & 1
\end{pmatrix}
\begin{pmatrix}
1 \\ 0 \\ 0
\end{pmatrix}
=
\begin{pmatrix}
1 \\ i \\ 0
\end{pmatrix}.
\]

\paragraph{Second Iteration:}
\[
v_2 = A v_1 = 
\begin{pmatrix}
1 & i & 0 \\
i & 1 & i \\
0 & i & 1
\end{pmatrix}
\begin{pmatrix}
1 \\ i \\ 0
\end{pmatrix}
=
\begin{pmatrix}
1 + i \\ 2i \\ i
\end{pmatrix}.
\]

\paragraph{Third Iteration:}
\[
v_3 = A v_2 = 
\begin{pmatrix}
1 & i & 0 \\
i & 1 & i \\
0 & i & 1
\end{pmatrix}
\begin{pmatrix}
1 + i \\ 2i \\ i
\end{pmatrix}
=
\begin{pmatrix}
0 \\ 0 \\ i
\end{pmatrix}.
\]

This sequence exhibits a Fibonacci-like recursion with complex coefficients, where the imaginary unit \(i\) introduces rotational and twisting effects.

\subsection*{Connections to Puboid Numbers}

Demi Moorish saw immediate parallels between the complex Fibonacci matrix and her puboid numbers:
\begin{itemize}
    \item \textbf{Matrix as a Puboid Representation:} The rows of \(A\) correspond to the three dimensions of a puboid, with the imaginary entries encoding rotations or twists in higher-dimensional space.
    \item \textbf{Determinant and Puboids:} The determinant of \(A\), \(3\), matches the triple structure of puboids formed by three distinct primes. This suggests a geometric framework for understanding puboid partitions.
    \item \textbf{Recursive Growth:} The Fibonacci-like recursion in \(A\) mirrors the growth of Bell numbers, hinting at deeper ties between partitioning and geometric symmetry.
\end{itemize}
Marjorie Worpitsky realized that her Pascalian polynomials could be extended to incorporate the structure of \(A\):
\begin{itemize}
    \item \textbf{Pascalian Polynomials with Complex Roots:} The recursive nature of Pascalian polynomials aligned with the iterations of the complex Fibonacci matrix. By introducing \(i\) into the coefficients, Marjorie proposed a new class of Pascalian polynomials tied to complex symmetries.
    \item \textbf{Geometric Interpretation:} The Pascalian roots of unity generalized to include rotations and twists in complex spaces, providing a richer algebraic structure for exploring symmetry.
\end{itemize}

The collaboration between Marjorie and Demi unfolded new questions:
\begin{itemize}
    \item \textbf{Eigenvalues and Stability:} What do the eigenvalues of \(A\) reveal about the behavior of the complex Fibonacci sequence? Are there ties to the roots of unity or modular symmetries?
    \item \textbf{Scaling and Partitions:} Can the determinant of \(A\) and its iterations model scaling laws or factorization properties of puboid numbers?
    \item \textbf{Higher Dimensions:} How might \(A\) be generalized to larger matrices, and what new symmetries could emerge from these extensions?
\end{itemize}

As Marjorie and Demi worked late into the night, the threads of their ideas began weaving into a larger tapestry. The complex Fibonacci matrix, Pascalian polynomials, and puboid numbers all pointed to a hidden harmony between recursion, geometry, and algebra.

\emph{"Perhaps,"} Marjorie mused, \emph{"our work isn’t just about structures or sequences—it’s about discovering the symmetries that underlie them all."}

With that, they joined the others in the Studio of Equilibrium, ready to share their progress and chart new directions for their explorations.

